% ==============================================================================
% Paper Template — Deepwork Research Platform
% Targeting NeurIPS 2026 style. Adjust \usepackage{neurips_2026} for other venues.
% ==============================================================================

\documentclass{article}

% --- Venue style ---
% Swap this line for your target venue's style file.
% Options: neurips_2026, icml2026, iclr2026_conference
\usepackage[preprint]{neurips_2026}

% --- Core packages ---
\usepackage[utf8]{inputenc}
\usepackage[T1]{fontenc}
\usepackage{hyperref}
\usepackage{url}
\usepackage{booktabs}       % Professional tables
\usepackage{amsfonts}
\usepackage{amsmath}
\usepackage{amsthm}
\usepackage{amssymb}
\usepackage{mathtools}
\usepackage{nicefrac}
\usepackage{microtype}       % Better typography
\usepackage{xcolor}

% --- Figures and graphics ---
\usepackage{graphicx}
\usepackage{subcaption}

% --- Algorithms ---
\usepackage[ruled,vlined,linesnumbered]{algorithm2e}

% --- Tables ---
\usepackage{multirow}
\usepackage{makecell}

% --- Theorem environments ---
\newtheorem{theorem}{Theorem}
\newtheorem{lemma}[theorem]{Lemma}
\newtheorem{proposition}[theorem]{Proposition}
\newtheorem{corollary}[theorem]{Corollary}
\newtheorem{conjecture}[theorem]{Conjecture}

\theoremstyle{definition}
\newtheorem{definition}[theorem]{Definition}
\newtheorem{example}[theorem]{Example}
\newtheorem{assumption}[theorem]{Assumption}

\theoremstyle{remark}
\newtheorem{remark}[theorem]{Remark}

% --- Common math macros ---
\newcommand{\R}{\mathbb{R}}
\newcommand{\N}{\mathbb{N}}
\newcommand{\Z}{\mathbb{Z}}
\newcommand{\Q}{\mathbb{Q}}
\newcommand{\C}{\mathbb{C}}
\newcommand{\E}{\mathbb{E}}
\newcommand{\Prob}{\mathbb{P}}
\newcommand{\Var}{\mathrm{Var}}
\newcommand{\Cov}{\mathrm{Cov}}
\DeclareMathOperator*{\argmax}{arg\,max}
\DeclareMathOperator*{\argmin}{arg\,min}
\newcommand{\ind}{\mathbbm{1}}
\newcommand{\norm}[1]{\left\| #1 \right\|}
\newcommand{\abs}[1]{\left| #1 \right|}
\newcommand{\inner}[2]{\langle #1, #2 \rangle}
\newcommand{\set}[1]{\left\{ #1 \right\}}
\newcommand{\ceil}[1]{\left\lceil #1 \right\rceil}
\newcommand{\floor}[1]{\left\lfloor #1 \right\rfloor}
\newcommand{\KL}{\mathrm{KL}}
\newcommand{\softmax}{\mathrm{softmax}}
\newcommand{\sigmoid}{\sigma}

% --- Editing macros (remove before submission) ---
\newcommand{\todo}[1]{\textcolor{red}{\textbf{TODO:} #1}}
\newcommand{\note}[1]{\textcolor{blue}{\textbf{Note:} #1}}

% ==============================================================================

\title{[Paper Title]}

% For anonymous submission, comment out authors and use:
% \author{Anonymous}

\author{
  [Author Name] \\
  [Affiliation] \\
  \texttt{[email]} \\
  % Add more authors with \And
}

\begin{document}

\maketitle

% ==============================================================================
\begin{abstract}
% 150-250 words. Four sentences minimum:
% 1. What is the problem and why does it matter?
% 2. What is your approach/contribution?
% 3. What are the key results?
% 4. What is the broader implication?

\todo{Write abstract last, after all other sections are complete.}
\end{abstract}

% ==============================================================================
\section{Introduction}
\label{sec:introduction}

% Structure (typically 4-6 paragraphs):
% 1. Opening: The broad problem area and why it matters (cite foundational work)
% 2. Specific gap: What's missing in current understanding? (cite recent work showing the gap)
% 3. This work: What we do — state contributions clearly
% 4. Key results: The headline findings, briefly
% 5. Implications: Why these results matter
% 6. Paper outline: "The remainder of this paper is organized as follows..."

% Clearly state contributions as a bulleted list:
% Our main contributions are:
% \begin{itemize}
%   \item We propose...
%   \item We prove/show...
%   \item We empirically demonstrate...
% \end{itemize}

\todo{Write introduction.}

% ==============================================================================
\section{Related Work}
\label{sec:related_work}

% Organize by theme, not chronologically. Common structure:
% \paragraph{[Theme 1].} Discuss the most relevant prior work stream...
% \paragraph{[Theme 2].} Another related area...
% \paragraph{[Theme 3].} Additional context...
%
% End each paragraph by clearly differentiating our work.
% Aim for 30+ citations. Be generous and fair to related work.

\todo{Write related work.}

% ==============================================================================
\section{Preliminaries}
\label{sec:preliminaries}

% Define notation, key concepts, and problem setup.
% Include formal definitions that later sections build on.
% Only include what the reader needs — move supplementary background to the appendix.

\todo{Write preliminaries.}

% ==============================================================================
\section{Method}
\label{sec:method}
% Alternative titles: "Approach", "Framework", "Theoretical Framework", "[Name of Your Method]"

% This is the core contribution section. Structure depends on the paper type:
%
% Theory paper:
%   - Formal definitions → Main theorem → Proof sketch → Implications
%   - Full proofs in appendix if they're long
%
% Systems/methods paper:
%   - Problem formulation → Algorithm/Architecture → Analysis
%
% Empirical paper:
%   - Experimental framework → Benchmark design → Evaluation protocol

\todo{Write method section.}

% ==============================================================================
\section{Experiments}
\label{sec:experiments}

% Standard structure:
% 5.1 Experimental Setup
%   - Datasets, models, baselines, hyperparameters, compute budget
%   - Evaluation metrics with formal definitions
% 5.2 Main Results
%   - Table or figure with headline results
%   - Statistical significance and confidence intervals
% 5.3 Ablation Studies
%   - What happens when you remove each component?
% 5.4 Analysis
%   - Qualitative examples, failure cases, scaling behavior, etc.

\subsection{Experimental Setup}
\label{sec:experimental_setup}

\todo{Describe setup.}

\subsection{Main Results}
\label{sec:main_results}

\todo{Present main results.}

\subsection{Ablation Studies}
\label{sec:ablations}

\todo{Present ablations.}

\subsection{Analysis}
\label{sec:analysis}

\todo{Deeper analysis.}

% ==============================================================================
\section{Discussion}
\label{sec:discussion}

% Address:
% - Limitations of the current work (be honest — reviewers respect this)
% - Broader implications of the findings
% - Connections to other fields or applications
% - Ethical considerations, if applicable
% - Future work directions (be specific, not vague)

\todo{Write discussion.}

% ==============================================================================
\section{Conclusion}
\label{sec:conclusion}

% 1-2 paragraphs. Restate the problem, summarize what was done, highlight key findings,
% and end with the broader significance. Do not introduce new information.

\todo{Write conclusion.}

% ==============================================================================
% Acknowledgments (hidden during anonymous review)
\begin{ack}
\todo{Add acknowledgments.}
\end{ack}

% ==============================================================================
\bibliographystyle{plainnat}
\bibliography{references}

% ==============================================================================
\newpage
\appendix

\section{Proofs}
\label{app:proofs}

% Full proofs of theorems stated in the main text.

\todo{Add proofs.}

\section{Additional Experiments}
\label{app:additional_experiments}

% Supplementary experimental results, extended tables, additional figures.

\todo{Add supplementary experiments.}

\section{Hyperparameter Details}
\label{app:hyperparameters}

% Full hyperparameter configurations for reproducibility.

\todo{Add hyperparameter details.}

\end{document}
